\subsection{CPython Memory Optimization Architecture (PEP 393)}

Python's \texttt{str} object is a language-level text sequence. CPython, the standard Python implementation, must still store that text in real memory. Since Python strings can contain anything from plain ASCII to emoji, using the same fixed-width storage for every string would waste memory in common cases.

PEP 393 introduced CPython's flexible string representation. The basic idea is simple: CPython chooses an internal character storage width according to the largest codepoint contained in the string.

\begin{verbatim}
texts = [
    "ABC",
    "ééé",
    "ăπ☃",
    "😀😀😀",
]

for text in texts:
    max_codepoint = max(ord(ch) for ch in text)
    print(f"text={text!r:<10} max=U+{max_codepoint:06X}")
\end{verbatim}

Possible output:

\begin{verbatim}
text='ABC'      max=U+000043
text='ééé'      max=U+0000E9
text='ăπ☃'      max=U+002603
text='😀😀😀'    max=U+01F600
\end{verbatim}

The important value is not the first character, nor the average character. The important value is the largest codepoint. That value determines the smallest internal storage width that can hold every component of the string.

\subsubsection{The Flexible String Representation Framework: Dynamic Character Data Width Selection Based on Maximum Codepoint Value}

CPython's flexible representation can be understood with three ordinary ranges:

\begin{itemize}
    \item codepoints up to \texttt{U+00FF} can fit in one byte per codepoint;
    \item codepoints up to \texttt{U+FFFF} can fit in two bytes per codepoint;
    \item larger valid Unicode codepoints need four bytes per codepoint.
\end{itemize}

This is not UTF-8. It is an internal CPython storage strategy for the canonical string data.

\begin{verbatim}
texts = [
    "ABC",
    "ééé",
    "ăπ☃",
    "😀😀😀",
]

for text in texts:
    max_nr = max(ord(ch) for ch in text)

    if max_nr <= 0xFF:
        width = 1
    elif max_nr <= 0xFFFF:
        width = 2
    else:
        width = 4

    print(f"text={text!r:<10} max=U+{max_nr:06X} likely_width={width}")
\end{verbatim}

Possible output:

\begin{verbatim}
text='ABC'      max=U+000043 likely_width=1
text='ééé'      max=U+0000E9 likely_width=1
text='ăπ☃'      max=U+002603 likely_width=2
text='😀😀😀'    max=U+01F600 likely_width=4
\end{verbatim}

This program does not directly read CPython's private C fields. Normal Python code does not expose the internal string kind as a public attribute. The program only computes the expected storage category from the visible Unicode codepoints.

The memory effect can be observed indirectly with \texttt{sys.getsizeof()}.

\begin{verbatim}
import sys
import platform

texts = [
    ("ascii", "ABC"),
    ("latin1", "ééé"),
    ("bmp", "ăπ☃"),
    ("non_bmp", "😀😀😀"),
]

print(f"implementation: {platform.python_implementation()}")
print()

for label, text in texts:
    max_nr = max(ord(ch) for ch in text)
    size = sys.getsizeof(text)

    print(f"{label:<8} text={text!r:<10} len={len(text):<2} "
          f"max=U+{max_nr:06X} size={size}")
\end{verbatim}

Possible output on one 64-bit CPython build:

\begin{verbatim}
implementation: CPython

ascii    text='ABC'      len=3  max=U+000043 size=44
latin1   text='ééé'      len=3  max=U+0000E9 size=60
bmp      text='ăπ☃'      len=3  max=U+002603 size=64
non_bmp  text='😀😀😀'    len=3  max=U+01F600 size=72
\end{verbatim}

The exact numbers are version-dependent and platform-dependent. The important pattern is that the internal representation is not simply ``one Python character equals one byte''.

Also, \texttt{sys.getsizeof()} reports the memory directly attributed to the object. It is useful for demonstrations, not a complete model of all memory activity in a running process.

The flexible representation also keeps \texttt{len()} simple. Length is still the number of Python string components, not the number of UTF-8 bytes.

\begin{verbatim}
texts = [
    "ABC",
    "ééé",
    "ăπ☃",
    "😀😀😀",
]

for text in texts:
    raw = text.encode("utf-8")
    print(f"text={text!r:<10} len(str)={len(text):<2} len(utf8)={len(raw):<2}")
\end{verbatim}

Possible output:

\begin{verbatim}
text='ABC'      len(str)=3  len(utf8)=3
text='ééé'      len(str)=3  len(utf8)=6
text='ăπ☃'      len(str)=3  len(utf8)=7
text='😀😀😀'    len(str)=3  len(utf8)=12
\end{verbatim}

The runtime string and its external UTF-8 byte serialization are different objects with different storage rules.

\subsubsection{Core Immutability: Fixed Heap Allocations and Read-Only Character Storage}

A CPython string object is immutable. Once the object exists, its character sequence cannot be edited in place from Python code.

\begin{verbatim}
text = "D'oh!"

print(f"text before: {text!r}")
print(f"id before:   {id(text)}")

try:
    text[0] = "W"
except TypeError as err:
    print("mutation failed")
    print(f"message: {err}")

print()
print(f"text after:  {text!r}")
print(f"id after:    {id(text)}")
\end{verbatim}

Possible output:

\begin{verbatim}
text before: "D'oh!"
id before:   140225874655600
mutation failed
message: 'str' object does not support item assignment

text after:  "D'oh!"
id after:    140225874655600
\end{verbatim}

The exact \texttt{id()} value is not important. The useful fact is that the existing string object was not edited.

String operations that appear to change text actually create another string object.

\begin{verbatim}
text = "spam"

old_id = id(text)
text = text.upper()

print(f"text:      {text!r}")
print(f"old id:    {old_id}")
print(f"new id:    {id(text)}")
print(f"same obj:  {old_id == id(text)}")
\end{verbatim}

Possible output:

\begin{verbatim}
text:      'SPAM'
old id:    140225874655920
new id:    140225874656176
same obj:  False
\end{verbatim}

The name \texttt{text} now refers to a different string object. The old object was not modified.

The same rule applies to concatenation.

\begin{verbatim}
text = "spam"

old_id = id(text)
text = text + " and eggs"

print(f"text:     {text!r}")
print(f"old id:   {old_id}")
print(f"new id:   {id(text)}")
print(f"same obj: {old_id == id(text)}")
\end{verbatim}

Possible output:

\begin{verbatim}
text:     'spam and eggs'
old id:   140225874655920
new id:   140225874644720
same obj: False
\end{verbatim}

This connects directly to the earlier distinction between rebinding and mutation. The assignment sign does not edit the old string. It binds the name to the result object.

Because strings are immutable, CPython can safely reuse, cache, hash, and intern some string objects. If a string could be changed in place, those optimizations would be dangerous.

\subsubsection{Computational and Allocation Overhead of Iterative String Concatenation}

Since strings are immutable, repeated concatenation can create many temporary string objects. Each new result must contain the previous text plus the newly added text.

The following program shows the identity changes during repeated concatenation.

\begin{verbatim}
text = ""

for part in ["Nobody", " expects", " the Spanish", " Inquisition"]:
    old_id = id(text)
    text = text + part

    print(f"added={part!r:<15} len={len(text):<3} "
          f"old_id={old_id} new_id={id(text)}")

print()
print(f"final text: {text!r}")
\end{verbatim}

Possible output:

\begin{verbatim}
added='Nobody'        len=6   old_id=140225874320944 new_id=140225874656112
added=' expects'      len=14  old_id=140225874656112 new_id=140225874644720
added=' the Spanish'  len=26  old_id=140225874644720 new_id=140225874646320
added=' Inquisition'  len=38  old_id=140225874646320 new_id=140225874646512

final text: 'Nobody expects the Spanish Inquisition'
\end{verbatim}

The repeated \texttt{+} pattern is easy to read for a small number of pieces. It is not the right construction pattern for many pieces.

The preferred pattern is:

\begin{enumerate}
    \item collect the pieces in a list;
    \item join them once at the end.
\end{enumerate}

\begin{verbatim}
parts = []

parts.append("Nobody")
parts.append(" expects")
parts.append(" the Spanish")
parts.append(" Inquisition")

text = "".join(parts)

print(f"parts: {parts!r}")
print(f"text:  {text!r}")
\end{verbatim}

Possible output:

\begin{verbatim}
parts: ['Nobody', ' expects', ' the Spanish', ' Inquisition']
text:  'Nobody expects the Spanish Inquisition'
\end{verbatim}

The list is mutable, so \texttt{append()} changes the list structure in place. The final \texttt{join()} operation computes the required final string once.

A larger demonstration can compare the two styles.

\begin{verbatim}
import time

count = 50000
piece = "ha"

start = time.perf_counter()

text = ""
for _ in range(count):
    text = text + piece

end = time.perf_counter()
concat_time = end - start

start = time.perf_counter()

parts = []
for _ in range(count):
    parts.append(piece)

text2 = "".join(parts)

end = time.perf_counter()
join_time = end - start

print(f"same result: {text == text2}")
print(f"len result:  {len(text)}")
print(f"+ loop:      {concat_time:.6f} seconds")
print(f"join style:  {join_time:.6f} seconds")
\end{verbatim}

Possible output:

\begin{verbatim}
same result: True
len result:  100000
+ loop:      0.090000 seconds
join style:  0.003000 seconds
\end{verbatim}

The exact timings depend on the machine, Python version, and current system load. The structural point is more important than the exact numbers: repeated immutable concatenation can allocate and copy repeatedly, while list accumulation followed by \texttt{join()} is the standard construction pattern.

For short direct expressions, \texttt{+} is fine.

\begin{verbatim}
name = "Kramer"
line = "Hello, " + name + "!"

print(line)
\end{verbatim}

Possible output:

\begin{verbatim}
Hello, Kramer!
\end{verbatim}

The warning is not ``never use \texttt{+} with strings''. The warning is: do not build long strings by repeated concatenation inside large loops when a list plus \texttt{join()} expresses the construction more directly.

\subsubsection{The Interning Subsystem: Immutable String Optimization and Singly Allocated Literals inside CPython}

Interning is an optimization where CPython keeps one shared string object for selected equal string values. If two names refer to the same interned string object, identity comparison with \texttt{is} will be true.

This is not the same as ordinary equality.

\begin{verbatim}
a = "python"
b = "python"

print(f"a == b: {a == b}")
print(f"a is b: {a is b}")
print(f"id(a):  {id(a)}")
print(f"id(b):  {id(b)}")
\end{verbatim}

Possible output on CPython:

\begin{verbatim}
a == b: True
a is b: True
id(a):  140225874658096
id(b):  140225874658096
\end{verbatim}

The result above is common for simple string literals, but program logic must not depend on it. Equality asks whether the text values are the same. Identity asks whether the two names refer to the same object.

A dynamically built string may be equal without being the same object.

\begin{verbatim}
a = "python"
b = "".join(["py", "thon"])

print(f"a: {a!r}")
print(f"b: {b!r}")

print()
print(f"a == b: {a == b}")
print(f"a is b: {a is b}")

print()
print(f"id(a): {id(a)}")
print(f"id(b): {id(b)}")
\end{verbatim}

Possible output:

\begin{verbatim}
a: 'python'
b: 'python'

a == b: True
a is b: False

id(a): 140225874658096
id(b): 140225874657840
\end{verbatim}

The text is equal. The objects may be different.

The \texttt{sys.intern()} function explicitly asks CPython to intern a string.

\begin{verbatim}
import sys

a = "python"
b = "".join(["py", "thon"])

print("before interning")
print(f"a == b: {a == b}")
print(f"a is b: {a is b}")

print()
a = sys.intern(a)
b = sys.intern(b)

print("after interning")
print(f"a == b: {a == b}")
print(f"a is b: {a is b}")

print()
print(f"id(a): {id(a)}")
print(f"id(b): {id(b)}")
\end{verbatim}

Possible output:

\begin{verbatim}
before interning
a == b: True
a is b: False

after interning
a == b: True
a is b: True

id(a): 140225874658096
id(b): 140225874658096
\end{verbatim}

Interning is useful for repeated string keys, especially identifier-like strings used in dictionaries. If both the stored key and the lookup key are interned, CPython can sometimes use pointer comparison after hashing instead of comparing the full character sequence.

For ordinary beginner-level code, the practical rule is still simple:

\begin{quote}
Use \texttt{==} to compare string text. Treat \texttt{is} results on strings as CPython optimization evidence, not as string equality logic.
\end{quote}

The following example shows why.

\begin{verbatim}
word1 = "springfield"
word2 = "spring" + "field"
word3 = "".join(["spring", "field"])

print(f"word1 == word2: {word1 == word2}")
print(f"word1 == word3: {word1 == word3}")

print()
print(f"word1 is word2: {word1 is word2}")
print(f"word1 is word3: {word1 is word3}")
\end{verbatim}

Possible output on CPython:

\begin{verbatim}
word1 == word2: True
word1 == word3: True

word1 is word2: True
word1 is word3: False
\end{verbatim}

The first two expressions may be folded or interned by CPython because their values are visible early. The third one is constructed at runtime. These details can vary. The semantic test remains \texttt{==}.

Interning is possible because strings are immutable. If the shared interned object could be edited in place, every name pointing to it would see the change, and dictionary key behavior would become unsafe. Immutable storage makes sharing safe.